\label{sec:static}
On the \val{NsnapTest} test snapshots, regret-2 repair finished within about one millisecond but left a regret of
\val{RepairRegret}\%. With a budget of 250\,ms, ALNS with the full budget reached \val{FixedRegret}\%, whereas
AnytimeSwarm reached \val{ASRegret}\% with \val{ASUsed}\% of the budget, that is, with \val{SavedPct}\% less computation.
Because the stopping rules use different amounts of computation, they are best compared at equal computation.
Table~\ref{tab:scale} makes this comparison for each fleet scale, with the fixed budget and the stagnation rule evaluated
at the mean computation that AnytimeSwarm used on the same scale. AnytimeSwarm reached a lower regret on every scale,
and all paired intervals lie below zero. The reduction was \val{ScMediumRedMin}\% to \val{ScMediumRedMax}\% on the
medium scale, \val{ScLargeRedMin}\% to \val{ScLargeRedMax}\% on the large scale and \val{ScXlargeRedMin}\% to
\val{ScXlargeRedMax}\% with 40 UAVs, a scale that was never used for training. Read in the other direction, the two
baselines needed between \val{ScMediumSpeedMin} and \val{ScSmallSpeedMax} times the computation of AnytimeSwarm to
reach its regret. Over all scales, AnytimeSwarm reached \val{EqRegretAS}\% against \val{EqRegretBaseMin}\% to
\val{EqRegretBaseMax}\% for four baselines that include a fixed budget and a stagnation patience tuned per fleet scale
on development data.

\inputtable{tab_static_scale}

Fig.~\ref{fig:static} shows where this advantage comes from. On the medium and large scales, the quality--compute
frontier of AnytimeSwarm lies below those of fixed budgets and stagnation stopping over the range of practical
computation (Fig.~\ref{fig:static}a). Grouped by the improvement that a 500\,ms search can still find
(Fig.~\ref{fig:static}b), AnytimeSwarm spent \val{AllocASNone}\,ms on snapshots that cannot be improved and
\val{AllocASHigh}\,ms on those with the largest available improvement. Only \val{AllocASShare}\% of its computation went
to snapshots without any possible gain, compared with more than half for the baselines. The learned predictor is
essential for this allocation. It reached an area under the ROC curve~\cite{fawcett2006roc} of \val{AUCLearned} for the
event that the next 10\,ms bring an improvement, compared with \val{AUCOnline} for an extrapolation of the recent
improvement rate.

\begin{figure}[t]
\centering
\begin{minipage}[b]{0.62\linewidth}\centering\fig{fig_static_frontier}{\linewidth}\\[-1pt]{\small (a)}\end{minipage}\hfill
\begin{minipage}[b]{0.36\linewidth}\centering\fig{fig_static_allocation}{\linewidth}\\[-1pt]{\small (b)}\end{minipage}
\caption{Exact execution, test seeds 2000--2029: (a) quality--compute frontiers on identical search trajectories (open
markers: baselines tuned per scale; star: selected AnytimeSwarm setting); (b) mean computation per snapshot, grouped by
the improvement that a 500\,ms search can still find.}\label{fig:static}
\end{figure}

\subsection{Reproduction on the evaluation platform}\label{sec:repro}
The exact-execution benchmark was produced on Windows. Before the execution-aware study, its development pipeline was
re-run on the platform of all later experiments (Ubuntu 24.04 under WSL2, identical Python and package versions). The
deterministic part was reproduced bit for bit, except for differences of one unit in the last place of some energy
values, which stem from the floating-point libraries of the two operating systems. The search ran
\val{JReproSpeedup}\% more iterations per millisecond on the new platform. The improvement predictor, re-trained on the
new timing traces, reached the same compute price by the rule fixed in the protocol, and the conclusions of the
equal-computation comparison agreed in sign and significance in \val{JReproMatched} comparisons. The frozen predictor
was used unchanged in all later experiments.
